\documentclass[10pt]{article}
\usepackage[margin=0.7in]{geometry}
\usepackage[T1]{fontenc}
\usepackage{lmodern}
\usepackage{enumitem}
\setlist[enumerate]{leftmargin=*, itemsep=1.2em}
\setlist[enumerate,2]{label=(\Alph*), leftmargin=2em, itemsep=0.05em, topsep=0.25em}
\pagestyle{plain}

\begin{document}

\begin{center}
  {\large\bfseries Video 0 Quiz: Course Introduction}\\[0.25em]
  \textit{Choose the best answer for each question.}
\end{center}

\noindent Name:\ \rule{3.2in}{0.4pt}\hfill Date:\ \rule{1.25in}{0.4pt}

\vspace{0.7em}

\begin{enumerate}
  \item Which preparation sequence does the video recommend before a class meeting?
  \begin{enumerate}
    \item Watch the matching video, read the matching notes, then try the homework problems.
    \item Complete the homework first, then watch the video only if a problem is difficult.
    \item Attend class without preparation and use the notes only before an exam.
    \item Read research papers first and postpone the videos until the end of the course.
  \end{enumerate}

  \item In the first example problem, why does the video introduce the error matrix $E$ in $\widehat{A}=A+E$?
  \begin{enumerate}
    \item To model uncertainty in the known matrix and study how it affects the computed solution.
    \item To guarantee that every matrix has an inverse.
    \item To replace vectors by scalar quantities.
    \item To show that the exact solution $x$ is always zero.
  \end{enumerate}

  \item What criterion is used in the second example to choose a polynomial approximation to a function?
  \begin{enumerate}
    \item Minimize the maximum value of the polynomial.
    \item Minimize the number of polynomial coefficients.
    \item Minimize the mean squared error between the function and the polynomial.
    \item Force the polynomial to pass through every point of the function.
  \end{enumerate}

  \item What lower-bound idea for a convex function is highlighted in the third example?
  \begin{enumerate}
    \item A convex function lies below every chord.
    \item A convex function lies above each of its tangent lines.
    \item Every local maximum of a convex function is global.
    \item A convex function must be a polynomial.
  \end{enumerate}
\end{enumerate}

\vspace{0.7em}
\noindent\textit{This quiz assesses the main ideas from the course introduction video.}

\end{document}
